\begin{abstract}
When does pretraining change an independent task's linear readability, and how large or persistent can that change be? For a convolution--square--pool encoder, local second moments give an exact representation for arbitrary input distributions. Standard population regression then connects multi-output tasks to subspace geometry. We separate exterior encoder updates caused by objective/data anisotropy from those caused by parameter mobility, obtaining quantitative near-conservation for linear factorization. In a solvable quadratic family with a genuine encoder bottleneck, filter correlation and head coupling create new second-order directions. An open cone of tasks shares a fixed improvement followed by excess risk above initialization, with explicit gaps, finite arrival times, and a structured five-percent covariance range. This family also identifies a symmetry under which learning-rate ratios cannot change the ordered subspace path. More generally, initial task cross-block mass and normal motion budget bound finite risk changes, while the two projector tangents characterize which tasks admit rate competition. Source-only interval bounds give nonzero finite-window predictions on existing matched-function trajectories; saved future risks satisfy those intervals at a common energy section. Earlier certified projection rebound and population/empirical texture results provide complementary evidence. CIFAR-10 and DTD remain boundary cases: pre-existing head-selection gaps and weak class gains have not yet been explained by measured values of the new bounds. The scope is task-conditioned population geometry, not a universal learning-rate rule or a finite-sample generalization guarantee.
\end{abstract}
